The comparative simulation results presented in Sections 4.4 and 4.5 provide computational evidence that the prevailing ``Status Quo'' strategy---equal distribution of limited municipal funds across all barangays---is a mechanism of dilution that prevents any single jurisdiction from reaching the critical intervention threshold required for behavioral change. In contrast, the Heuristic-Guided Deep Reinforcement Learning (HuDRL) agent independently discovered a ``Sequential Saturation'' strategy: concentrating resources to force localized compliance past the simulated 70\% social tipping point, locking in behavior via the Social Norm Shield, and then rotating resources to the next target. This Section translates these computational findings into actionable policy insights for the Municipality of Bacolod, organized according to the four pillars of the proposed Data-Driven Solid Waste Management (DD-SWM) Framework.

\subsection{Algorithmic Resource Targeting (Sequential Saturation)}

Section 4.4.2 demonstrated mathematically that the Status Quo strategy of equal distribution is structurally insufficient. Even after 12 quarters of sustained funding, the Global Compliance Rate under Status Quo plateaued at 30.18\%, with the massive urban center of Barangay Poblacion ($N=1,534$) flatlining at 0.06\% compliance, while the smaller populations of Mati ($N=165$) and Babalaya ($N=171$) surged past 90\% with the exact same per-barangay allocation. This bifurcation reveals a critical policy failure: equal distribution produces neither equity nor efficacy. In contrast, the HuDRL agent (Section 4.4.5) concentrated municipal funds asymmetrically, directing 50\% of the quarterly budget to Poblacion in Quarter 1, forcing its compliance to 90.09\% by Quarter 2, and then rotating to rural barangays, achieving a terminal municipal-wide compliance of 90.45\% without exceeding the fixed \textpeso 375,000 quarterly budget.

The Municipal Mayor and the Sangguniang Bayan must abandon static, equal-share appropriation ordinances in favor of a legally ring-fenced “Sequential Saturation” strategy. Rather than distributing the \textpeso 1.5 million annual augmentation fund equally across all seven barangays, the LGU should adopt a Phased Rollout formally approved by the Municipal Solid Waste Management Board. The Mayor and the MENRO should publicly declare a prioritized “Pilot Zone” (e.g., Barangay Poblacion) as the initial recipient of 50--70\% of the quarterly budget for a fixed duration---most likely one quarter, based on the simulation---before rotating to the next target. After a barangay crosses 70\% compliance, its funding must be drastically titrated to a fractional “Maintenance Budget” (ideally 5--10\% of the total quarterly fund) to free capital for the next target.

Real-world implementation confronts the political reality of Philippine local governance. Elected officials are acutely sensitive to accusations of bias and patronage, particularly when resources are visibly concentrated on one barangay. Research demonstrates the “paradox of prioritization,” where algorithmic tools can intensify public perceptions of inequality \citep{Alphonso2024}. To depoliticize this unequal budget distribution, the LGU can use the simulation results as a “mathematical shield.” The Sangguniang Bayan should pass a resolution explicitly adopting the Sequential Saturation framework, accompanied by transparent, publicized criteria for barangay prioritization (e.g., baseline compliance, logistical accessibility), a fixed rotation schedule, and independent monitoring by credible civil society organizations. This transforms the Sequential Saturation strategy from a politically vulnerable executive decision into a codified, transparent governance framework.

\noindent \textbf{Computational optimality versus political reality.}

While the Sequential Saturation strategy is mathematically optimal within the simulation, its real-world feasibility is severely constrained by the political economy of Philippine local governance. The model assumes a rational, public-interest maximizing Mayor Agent; however, the actual allocation of LGU budgets is heavily shaped by patronage politics, pork barrel logics, and the prevalence of political dynasties. Former Supreme Court Justice Carpio has warned that entrenched political dynasties control approximately 80\% of Congress and 50\% of LGUs, using pork barrel allocations—constituting about 60\% of the annual budget—to ensure electoral victories for family members \citep{Carpio2025a}. Dynasties “often operate with little transparency and minimal accountability,” and when power is held for a long time by one family, “public policies begin to serve private, not public, interests” \citep{Tribune2025}. 

At the national level, an estimated half a trillion pesos remain tied to pork-risky road and other infrastructure projects, while the 2026 budget increased local government pork allocations to Php57.9 billion \citep{IBON2026}. Under such conditions, an algorithmically derived targeting scheme is unlikely to be implemented faithfully; rather, the discretion embedded in “transparent prioritization criteria” would be vulnerable to political capture, turning the Sequential Saturation framework from a neutral optimization tool into another instrument for pork barrel or a weapon to punish political opponents \citep{PIDS2024}. The PIDS has documented that even with increased fiscal autonomy under the Mandanas ruling, LGUs often struggle to align spending with devolved functions, engaging in “cross-charging” that obscures spending and undermines accountability \citep{PIDS2024}. Moreover, Philippine LGUs face binding constraints that include “fiscal imbalances, institutional capacity, enforcement, and political economy factors,” resulting in significant gaps between de jure and de facto functional assignments \citep{PIDS2025a}. 

Given these structural realities, the proposed Sequential Saturation policy is computationally optimal but politically infeasible under current governance conditions. The model’s output should therefore be interpreted not as a directly implementable blueprint, but as a \textit{proof of concept} that highlights the gap between algorithmic efficiency and on-the-ground political economy. Any attempt to pilot the strategy would require prior institutional reforms—such as civil society oversight, anti-dynasty legislation, and strengthened local accountability mechanisms—that lie well beyond the scope of this thesis \citep{Carpio2025a,IBON2026,Tribune2025,PIDS2024,PIDS2025a}.

\subsection{Friction Reduction and Infrastructure}
The Global Sensitivity Analysis (Section 4.2) revealed that the Cost of Effort ($C_{Effort}$) is the dominant driver of compliance variance within the model ($S_T=1.06$). In contrast, the Attitude weight exhibited near-zero first-order sensitivity. This provides strong computational evidence of a ``Value-Action Gap'': citizens are not failing to segregate because they lack awareness; they are actively choosing not to comply because the physical hassle remains consistently too high \citep{Yazawa2025, Zhao2022}.

The MENRO must restructure its budget allocation, shifting resources from IEC campaigns toward direct subsidies that reduce physical friction. The LGU should:
\begin{itemize}
    \item \textbf{Subsidize segregation receptacles} for low-income households, directly reducing $C_{Effort}$ at the point of generation.
    \item \textbf{Prioritize hardware maintenance}, ensuring collection vehicles are operational so the ``No Segregation, No Collection'' rule is logistically feasible.
    \item \textbf{Institutionalize purok-level collection with integrated enforcement.} Building upon the successful grassroots practices observed in Barangay Binuni, the LGU should implement a scaled-up, purok-based collection model. While Barangay Binuni currently utilizes a single plantilla collector equipped with a pushcart (kariton) who effectively collects waste and logs non-compliance at the barangay level (Appendix B.3), this concept can be optimized for broader municipal application. Under this proposed model, the LGU would hire garbage collectors at the purok level equipped with motorized cabs. These collectors would gather segregated waste directly from households, eliminating the need for residents to walk to a central MRF—a direct reduction of $C_{Effort}$. Crucially, the collector also serves as a waste inspector. When a household fails to segregate, the collector still accepts the waste but issues a fine and logs the violation, bringing the case to the barangay. The unsorted waste must then be segregated at the MRF, imposing extra work on the collector. This creates a powerful deterrence effect: collectors are personally motivated to identify and report non‑compliant households because such households directly generate additional manual labor for them. This aligns with the principle of combining monitoring with immediate disincentives \citep{Ma2023}, where visible, on‑the‑ground enforcement—not abstract regulations—changes behavior at the point of collection. This decentralized model has been codified in other Philippine LGUs; Lapu‑Lapu City Ordinance No. 17‑016‑2025, for example, explicitly institutionalizes a purok‑based SWM system that places responsibility for segregation and collection at the grassroots level, with barangays tasked to monitor compliance and coordinate with the CENRO. To sustain collector motivation and prevent corruption, collectors should be appointed as plantilla personnel with a regular salary, and quarterly rotation of assignment across puroks should be mandated to prevent familiarity and bias from undermining enforcement \citep{Ma2023}.
\end{itemize}

\subsection{Pillar III: Pulsed Enforcement, the Graduation Rule, and the Transactional Trap of Incentives}

A critical lesson from the Pure Incentives scenario is the danger of creating a purely transactional relationship. When the fixed municipal allocation was rapidly drained by early adopters, latecomers were left without compensation, triggering a ``Perceived Unfairness'' penalty. This aligns with the well-established finding that monetary incentives can undermine intrinsic motivation when perceived as controlling \citep{Chen2023, Zhao2022}.

The simulation reflects a well-documented governance phenomenon known as the ``Cobra Effect,'' where well-intentioned schemes produce unintended negative consequences \citep{Weaver1986}. The LGU must treat material incentives as a catalytic, time-limited tool, not as a permanent subsidy. Any incentive program must be explicitly designed with a fixed duration and a clear phase-out schedule that transitions households from reward-driven to norm-driven compliance.

The MENRO must adopt an adaptive enforcement strategy:
\begin{itemize}
    \item \textbf{Phase 1 (Saturation Phase):} Deploy the Municipal Strike Force intensively and activate temporary incentives to forcibly raise compliance past the tipping point.
    \item \textbf{Phase 2 (Norm Lock-in):} Once compliance exceeds 70\%, immediately scale back enforcement and incentives to a fractional maintenance level (The Graduation Rule) to prevent eroding intrinsic motivation.
    \item \textbf{Phase 3 (Rotation):} Redirect freed capital to the next target barangay.
\end{itemize}

\subsection{Pillar IV: Predictive Policy Prototyping}
The coupled ABM-DRL framework serves as a proof-of-concept for predictive policy prototyping. The HuDRL agent successfully explored a 21-dimensional continuous action space and discovered an adaptive strategy that outperformed all static heuristics \citep{Jimenez2025, ZhengAI2022}.

The Municipal Government should commission a customized decision-support system based on this DRL architecture. This tool can simulate the outcomes of appropriation ordinances before enactment, stress-test policies against logistical shocks, and provide quarterly policy recommendations to the MENRO. Given current LGU AI readiness levels, implementation should begin as a pilot in two to three barangays in partnership with the Department of Science and Technology (DOST), gradually expanding to cover the entire municipality. By providing a rigorous, data-driven baseline, the model shifts the LGU toward proactive, evidence-based governance.

\subsection{Synthesis: From Simulation to Governance Reality}
The four pillars constitute an integrated framework for transitioning the Municipality of Bacolod from reactive, heuristic-driven management to proactive, data-driven governance. However, three limitations must be transparently acknowledged: political feasibility remains a binding constraint necessitating transparent communication; infrastructure investment to reduce the Cost of Effort is a prerequisite, not an option; and the model's predictions remain conditional on its assumptions, serving as a testable hypothesis rather than a final blueprint.

Nevertheless, the comparative results within the simulation are unambiguous: equal distribution fails; single-instrument regimes fail; sequential saturation succeeds \citep{Alphonso2024, Jimenez2025}. For a resource-constrained LGU, the evidence suggests that abandoning the administrative heuristic of ``fairness through equality'' in favor of ``efficacy through sequential equity'' is the most promising path toward sustainable solid waste management.